% Generated by scripts/assemble.py; corrections belong in source/corrections.json.
\examhead{151}{国防科技大学 2026 年数学分析与高等代数试卷(回忆版)}{143}
% p144-b19; source PDF page 144
\blockanchor{p144-b19}
(注. 题目来自考生回忆, 部分表述或数据可能有误)\par
% p144-b20; source PDF page 144
\blockanchor{p144-b20}
\question{1}令 $D(x) = \left\{ \begin{array}{ll} 1, & x \text{为有理数}, \\ 0, & x \text{为无理数}. \end{array} \right.$\par
% p144-b21; source PDF page 144
\blockanchor{p144-b21}
(1) 证明: $D(x)$ 在实数范围内处处不存在极限.\par
% p144-b22; source PDF page 144
\blockanchor{p144-b22}
(2) 构造 $f(x)$, 使其在 $x = 1$ 和 $x = 2$ 处连续, 其余点处均不连续.\par
% p144-b23; source PDF page 144
\blockanchor{p144-b23}
\question{2}已知 $A^3 = 2E, B = A^2 - 2A + 2E,$ 求 $B^{-1}$.\par
% p145-b00; source PDF page 145
\par\Needspace{6\baselineskip}
\blockanchor{p145-b00}
\question{3}给定 $P^3$ 的两组基\par
\nopagebreak[4]
% p145-b01; source PDF page 145
\blockanchor{p145-b01}
\displayformula{\alpha_ {1} = (1, 0, 1), \alpha_ {2} = (2, 1, 0), \alpha_ {3} = (1, 1, 1);}
% p145-b02; source PDF page 145
\blockanchor{p145-b02}
\displayformula{\beta_ {1} = (1, 2, - 1), \beta_ {2} = (2, 2, - 1), \beta_ {3} = (2, - 1, - 1).}
% p145-b03; source PDF page 145
\blockanchor{p145-b03}
线性变换  $\sigma$  满足  $\sigma(\alpha_{i})=\beta_{i}(i=1,2,3)$ .\par
% p145-b04; source PDF page 145
\blockanchor{p145-b04}
(1) 求由基 $\alpha_{1}, \alpha_{2}, \alpha_{3}$ 到 $\beta_{1}, \beta_{2}, \beta_{3}$ 的过渡矩阵.\par
% p145-b05; source PDF page 145
\blockanchor{p145-b05}
(2) 写出 $\sigma$ 在基 $\alpha_{1}, \alpha_{2}, \alpha_{3}$ 下的矩阵.\par
% p145-b06; source PDF page 145
\blockanchor{p145-b06}
\question{4}未知.\par
% p145-b07; source PDF page 145
\blockanchor{p145-b07}
\question{5}计算 $\iint_{D} |(x - 1)^2 + (y - 1)^2 - 1| \, \mathrm{d}x \, \mathrm{d}y$, 其中 $D = \{(x, y) \mid x^2 + y^2 \leq 9\}$.\par
% p145-b08; source PDF page 145
\blockanchor{p145-b08}
\question{6}令 $A = \left( \begin{array}{ccc}1 & -3 & 1\\ 2 & 1 & 0\\ 3 & 1 & 1 \end{array} \right).$\par
% p145-b09; source PDF page 145
\blockanchor{p145-b09}
(1) 证明: $A$ 在有理数域上不可以对角化.\par
% p145-b10; source PDF page 145
\blockanchor{p145-b10}
(2) 证明: $A$ 在复数域上可以对角化.\par
% p145-b11; source PDF page 145
\blockanchor{p145-b11}
\question{7}设 $f(x) = x^{m}$ ($m$ 为大于 2 的正整数), 证明: $f(x)$ 在 $[0, +\infty)$ 上不一致连续.\par
% p145-b12; source PDF page 145
\blockanchor{p145-b12}
\question{8}设 $f(x)$ 在 $[0,1]$ 上连续可导，且 $f(0) = 0$，证明：$\int_0^1 |f(x)f'(x)| \mathrm{d}x \leq \frac{1}{2} \int_0^1 [f'(x)]^2 \mathrm{d}x.$\par
% p145-b13; source PDF page 145
\blockanchor{p145-b13}
\question{9}证明: $a_{n} = \sum_{k=1}^{n} \frac{1}{n+k}$ 收敛.\par
% p145-b14; source PDF page 145
\par\Needspace{6\baselineskip}
\blockanchor{p145-b14}
\question{10}设 n 阶实对称矩阵 A 的全部特征值按大小顺序排列成  $\lambda_{1} \geq \lambda_{2} \geq \cdots \geq \lambda_{n}$ ，证明：\par
\nopagebreak[4]
% p145-b15; source PDF page 145
\blockanchor{p145-b15}
\displayformula{\max _ {x \neq 0} \frac {x ^ {\mathrm{T}} A x}{x ^ {\mathrm{T}} x} = \lambda_ {1}, \min _ {x \neq 0} \frac {x ^ {\mathrm{T}} A x}{x ^ {\mathrm{T}} x} = \lambda_ {n}.}
